% Archive record 52, line 680.
\hypertarget{part-4-52-0680-iv-a-ux521dux7a3fux4e2dux7684ux8865ux5145ux63a8ux5bfcux4e0eux4f8bux56fe}{%
\section{IV-A --- 初稿中的补充推导与例图}\label{part-4-52-0680-iv-a-ux521dux7a3fux4e2dux7684ux8865ux5145ux63a8ux5bfcux4e0eux4f8bux56fe}}

\hypertarget{part-4-52-0680-volatility-clusteringux6ce2ux52a8ux7387ux805aux96c6}{%
\subsubsection{6.2 Volatility Clustering（波动率聚集）}\label{part-4-52-0680-volatility-clusteringux6ce2ux52a8ux7387ux805aux96c6}}

金融市场中的高波动通常倾向于连续出现，低波动时期也可能持续一段时间。

\textbf{Volatility Clustering}

\bookgraphic{0.455}{part-4-52-0680-001-01}{合成示意：下方序列的收益率方向仍然难以预测，但波动幅度具有明显的持续性。}

因此，我们可能同时观察到：
\[\operatorname{Corr}(r_t,r_{t+h})\approx0,\]

但：
\[\operatorname{Corr}(|r_t|,|r_{t+h}|)>0,\]

或者：
\[\operatorname{Corr}(r_t^2,r_{t+h}^2)>0.\]

这意味着：

\textbf{收益率本身可能不具有明显的线性可预测性，但其波动程度具有可预测性。}

\hypertarget{part-4-52-0680-ux4e00ux4e2aux4e25ux683cux7684ux6982ux7387ux53cdux4f8b}{%
\subsubsection{6.3 一个严格的概率反例}\label{part-4-52-0680-ux4e00ux4e2aux4e25ux683cux7684ux6982ux7387ux53cdux4f8b}}

为了进一步理解这一点，我们构造一个非常简洁的随机过程。 设：
\[Z_t\overset{\mathrm{i.i.d.}}{\sim}\mathcal N(0,1),\]

并定义：
\[r_t=Z_tZ_{t-1}.\]

首先：
\[\mathbb E[r_t]=0.\]

相邻收益率的协方差为：
\[\begin{aligned}
\operatorname{Cov}(r_t,r_{t+1})
&=\mathbb E[Z_t^2Z_{t-1}Z_{t+1}]\\
&=\mathbb E[Z_t^2]
\mathbb E[Z_{t-1}]
\mathbb E[Z_{t+1}]\\
&=0.
\end{aligned}\]

但如果计算平方收益率：
\[r_t^2=Z_t^2Z_{t-1}^2,\]

由于标准正态分布满足：
\[\mathbb E[Z_t^2]=1,\qquad
\mathbb E[Z_t^4]=3,\]

可以得到：
\[\begin{aligned}
\operatorname{Cov}(r_t^2,r_{t+1}^2)
&=\mathbb E[Z_t^4Z_{t-1}^2Z_{t+1}^2]-1\\
&=3-1\\
&=2.
\end{aligned}\]

因此：
\[\boxed{
\begin{aligned}
\operatorname{Cov}(r_t,r_{t+1})&=0,\\
\operatorname{Cov}(r_t^2,r_{t+1}^2)&=2.
\end{aligned}
}\]

这严格证明了一个随机过程可以不存在收益率的线性相关性，却存在波动幅度的统计依赖。

当然，这只是一个教学反例，不是完整的金融市场模型。

\hypertarget{part-4-52-0680-arch-and-garch}{%
\subsubsection{6.4 ARCH and GARCH}\label{part-4-52-0680-arch-and-garch}}

为了描述 Volatility Clustering，金融计量经济学发展出了 ARCH 和 GARCH 模型。

目前只需要掌握 GARCH(1,1) 的基本结构。

\[r_t=\mu+\varepsilon_t,\]

\[\varepsilon_t=\sigma_tZ_t,\]

其中 \(Z_t\) 是均值为 0、方差为 1 的 IID Innovation。

条件方差由下式决定：
\[\boxed{
\sigma_t^2
=
\omega
+a\varepsilon_{t-1}^2
+b\sigma_{t-1}^2
}\]

这里新定义的参数满足：
\[\omega>0,\qquad a,b\geq0.\]

\begin{readertable}
\begin{longtable}[]{@{}
  >{\raggedright\arraybackslash}p{(\columnwidth - 2\tabcolsep) * \real{0.32}}
  >{\raggedright\arraybackslash}p{(\columnwidth - 2\tabcolsep) * \real{0.64}}@{}}
\toprule
Parameter & Meaning \\
\midrule
\endhead
\(\omega\) & 基础波动水平 \\
\(a\) & 对近期市场冲击的敏感程度 \\
\(b\) & 历史波动率的持续程度 \\
\(a+b\) & 波动冲击的持续性 \\
\bottomrule
\end{longtable}
\end{readertable}

在具有有限二阶矩的平稳解条件下，若 \(a+b<1\)，对条件方差方程两侧取期望：
\[\mathbb E[\sigma_t^2]
=
\omega
+(a+b)\mathbb E[\sigma_t^2].\]

于是：
\[\boxed{
\mathbb E[\sigma_t^2]
=
\frac{\omega}{1-a-b}
}\]

\textbf{GARCH --- Volatility Shock Persistence}

Persistence \(a+b\)\\
\textbf{0.90}

0.9（范围 0--0.99）

\bookgraphic{0.523}{part-4-52-0680-002-01}{}

Shock half-life

\textbf{6.58}

Periods

Remaining shock after 20 periods

\textbf{12.2\%}

归一化条件方差预测示意：长期方差为 1，初始超额冲击为 3。曲线描述预期方差的恢复过程，而非实际市场轨迹。

persistence = 0

\bookgraphic{0.523}{part-4-52-0680-002-state-01}{}

这里还有一个相当 Nontrivial 的观点：

\textbf{Conditional Variance 随时间变化，并不意味着 Unconditional Variance 一定随时间变化。}

GARCH 可以同时具有：

\begin{itemize}
\tightlist
\item
  Time-varying Conditional Variance；
\item
  Constant Unconditional Variance；
\item
  Weak Stationarity。
\end{itemize}

因此，Volatility Clustering 本身并不足以证明 Non-stationarity。

\hypertarget{part-4-52-0680-heavy-tailsux539aux5c3eux5206ux5e03}{%
\subsection{7. Heavy Tails（厚尾分布）}\label{part-4-52-0680-heavy-tailsux539aux5c3eux5206ux5e03}}

\hypertarget{part-4-52-0680-gaussian-assumption-ux7684ux5c40ux9650}{%
\subsubsection{7.1 Gaussian Assumption 的局限}\label{part-4-52-0680-gaussian-assumption-ux7684ux5c40ux9650}}

我们此前大量使用 Normal Distribution，因为它具有非常好的数学性质。

但真实金融收益率往往比正态分布产生更多极端观测值，这也是经典金融 Stylized Facts 之一。（\href{https://doi.org/10.1080/713665670}{Cont, 2001}）

\bookgraphic{0.473}{part-4-52-0680-004-01}{}

Standard Normal

Standardized Student's t (ν=3)

两个分布具有相同均值和方差，但 Student's t 在极端区域具有更高概率密度。

例如，将 Standard Normal Distribution 与标准化为单位方差的 Student's t Distribution（自由度为 3）进行比较。

Normal：超过 5 个标准差

\textbf{0.0000573\%}

Student's t：超过 5 个标准差

\textbf{0.324\%}

双侧尾部概率，基于理论分布计算，并非实际金融市场的实测频率。

即使均值和方差完全相同，两者的 Extreme Event Probability 也可能存在数量级差异。

\hypertarget{part-4-52-0680-moment-existence}{%
\subsubsection{7.2 Moment Existence}\label{part-4-52-0680-moment-existence}}

Student's t Distribution 还有一个值得关注的数学性质。

对于：
\[X\sim t_\nu,\]

其 \(p\) 阶绝对矩存在的条件为：
\[\boxed{
\mathbb E[|X|^p]<\infty
\iff p<\nu,\qquad p>0.
}\]

\begin{readertable}
\begin{longtable}[]{@{}
  >{\raggedright\arraybackslash}p{(\columnwidth - 6\tabcolsep) * \real{0.33}}
  >{\raggedright\arraybackslash}p{(\columnwidth - 6\tabcolsep) * \real{0.16}}
  >{\raggedright\arraybackslash}p{(\columnwidth - 6\tabcolsep) * \real{0.20}}
  >{\raggedright\arraybackslash}p{(\columnwidth - 6\tabcolsep) * \real{0.27}}@{}}
\toprule
Degrees of Freedom & Mean & Variance & Fourth Moment \\
\midrule
\endhead
\(\nu=1\) & 不存在 & 不存在 & 不存在 \\
\(\nu=2\) & 存在 & 不存在 & 不存在 \\
\(\nu=3\) & 存在 & 存在 & 不存在 \\
\(\nu=5\) & 存在 & 存在 & 存在 \\
\bottomrule
\end{longtable}
\end{readertable}

其中 \(t_1\) 就是你此前在 Probability 中学习过的 Cauchy Distribution。

这与 Part II 的统计理论产生直接联系：

如果某个 DGP 不满足有限方差或者有界损失条件，那么此前使用的 Classical CLT、Chebyshev Inequality、Hoeffding Inequality 等工具就不能无条件应用。

但也不能仅仅因为金融数据具有厚尾，就断言它一定没有有限方差。具体矩是否存在，取决于真实尾部分布。

\hypertarget{part-4-52-0680-temporal-dependence-ux4e0e-effective-sample-size}{%
\subsection{8. Temporal Dependence 与 Effective Sample Size}\label{part-4-52-0680-temporal-dependence-ux4e0e-effective-sample-size}}

\hypertarget{part-4-52-0680-ux4e2aux89c2ux6d4bux503cux4e0dux4e00ux5b9aux5bf9ux5e94-1000-ux4efdux72ecux7acbux4fe1ux606f}{%
\subsubsection{8.1 1,000 个观测值不一定对应 1,000 份独立信息}\label{part-4-52-0680-ux4e2aux89c2ux6d4bux503cux4e0dux4e00ux5b9aux5bf9ux5e94-1000-ux4efdux72ecux7acbux4fe1ux606f}}

此前 II-F 已经推导过，对于弱平稳序列：
\[\operatorname{Var}(\bar X_n)
=
\frac{\gamma_0}{n}
\left[
1+
2\sum_{h=1}^{n-1}
\left(1-\frac hn\right)\rho_h
\right].\]

如果所有非零阶自相关均为零，则恢复熟悉的：
\[\operatorname{Var}(\bar X_n)=\frac{\gamma_0}{n}.\]

为了刻画时间依赖对样本均值估计的影响，可以定义 Effective Sample Size：
\[\boxed{
n_{\mathrm{eff}}
=
\frac{n}{
1+
2\sum_{h=1}^{n-1}
(1-h/n)\rho_h
}
}\]

使得：
\[\operatorname{Var}(\bar X_n)
=
\frac{\gamma_0}{n_{\mathrm{eff}}}.\]

注意，这里的 Effective Sample Size 是针对样本均值方差定义的，不是所有统计量通用的有效样本数量。

\hypertarget{part-4-52-0680-ar1-example}{%
\subsubsection{8.2 AR(1) Example}\label{part-4-52-0680-ar1-example}}

如果：
\[\rho_h=\phi^h,\]

当 \(n\) 足够大且 \(|\phi|<1\) 时：
\[\begin{aligned}
1+2\sum_{h=1}^{\infty}\phi^h
&=1+\frac{2\phi}{1-\phi}\\
&=\frac{1+\phi}{1-\phi}.
\end{aligned}\]

因此：
\[\boxed{
n_{\mathrm{eff}}
\approx
n\frac{1-\phi}{1+\phi}
}\]

\textbf{Effective Sample Size Explorer}

Observed sample size \(n\)\\
\textbf{1,000}

1000（范围 100--5000）

Positive autocorrelation parameter \(\phi\)\\
\textbf{0.80}

0.8（范围 0--0.95）

Nominal sample size

\begin{verbatim}
  **1,000**
\end{verbatim}

Effective sample size

\begin{verbatim}
  **112**
\end{verbatim}

Effective information ratio

\textbf{11.2\%}

Standard error multiplier

\textbf{2.99×}

基于平稳 AR(1) 的有限样本均值方差公式计算。不是对真实金融数据有效样本量的估计。

这提供了一个非常有价值的 Research Insight。

在金融任务中，即使原始收益率的线性自相关很弱，模型损失、策略收益、波动率以及重叠预测标签仍然可能具有显著的时间依赖性。

因此，我们不能直接将原始数据的行数视为独立统计信息的数量。

\hypertarget{part-4-52-0680-ux4e3aux4ec0ux4e48ux91d1ux878dux73afux5883ux7279ux522bux5bb9ux6613ux4ea7ux751f-rsi-overfitting}{%
\subsection{9. 为什么金融环境特别容易产生 RSI Overfitting？}\label{part-4-52-0680-ux4e3aux4ec0ux4e48ux91d1ux878dux73afux5883ux7279ux522bux5bb9ux6613ux4ea7ux751f-rsi-overfitting}}

现在可以将本阶段的内容与此前的 Part II、Part III 连接起来。

\begin{readertable}
\begin{longtable}[]{@{}
  >{\raggedright\arraybackslash}p{(\columnwidth - 4\tabcolsep) * \real{0.29}}
  >{\raggedright\arraybackslash}p{(\columnwidth - 4\tabcolsep) * \real{0.31}}
  >{\raggedright\arraybackslash}p{(\columnwidth - 4\tabcolsep) * \real{0.37}}@{}}
\toprule
Financial Property & Statistical Consequence & RSI Consequence \\
\midrule
\endhead
Weak Predictability & 真实模型差异可能很小 & 容易选择噪声驱动的修改 \\
Temporal Dependence & 有效样本信息减少 & 评估可靠性下降 \\
Heavy Tails & 极端观测可能主导指标 & Agent 可能过度优化偶然事件 \\
Volatility Clustering & 评估误差具有条件异方差 & 不同市场窗口的评分稳定性不同 \\
Non-stationarity & 历史规律可能失效 & 后期版本可能过度适应旧环境 \\
Single Historical Path & 缺少独立市场历史 & 反复回测可能累积选择偏差 \\
\bottomrule
\end{longtable}
\end{readertable}

\textbf{Coding Tasks}

相对明确的正确性反馈

部分指标可确定性验证

任务之间可能存在迁移性

\textbf{Financial Tasks}

高统计不确定性

有限且相关的历史数据

潜在的环境变化

↓

\textbf{Same RSI Architecture} \(A_0\rightarrow A_1\rightarrow\cdots\rightarrow A_K\)

↓

\bookgraphic{0.629}{part-4-52-0680-008-01}{}

IS performance

Potential OOS trajectory in a reliable evaluation environment

Potential OOS trajectory in a low-SNR financial environment

研究假设示意，并非真实实验结果，也不意味着编程任务一定具有更好的泛化表现。

这里最重要的认识是：\textbf{金融环境未必削弱 Agent 的 Optimization Capability，却可能显著削弱 Evaluator 区分真实改进与虚假改进的能力。}

Agent 仍然可以变得越来越擅长搜索，但它所搜索到的高分候选未必代表真正的能力提升。

\hypertarget{part-4-52-0680-research-practice-initial-financial-data-audit}{%
\subsection{10. Research Practice --- Initial Financial Data Audit}\label{part-4-52-0680-research-practice-initial-financial-data-audit}}

如果未来课题组提供了一份金融数据，我们不应该立即开始运行 RSI。

第一步应当对 DGP 的统计特征进行初步分析。

\begin{readertable}
\begin{longtable}[]{@{}
  >{\raggedright\arraybackslash}p{(\columnwidth - 4\tabcolsep) * \real{0.32}}
  >{\raggedright\arraybackslash}p{(\columnwidth - 4\tabcolsep) * \real{0.37}}
  >{\raggedright\arraybackslash}p{(\columnwidth - 4\tabcolsep) * \real{0.27}}@{}}
\toprule
Diagnostic & Method & Research Purpose \\
\midrule
\endhead
Return Transformation & Simple / Log Return & 统一统计对象 \\
Distribution Analysis & Histogram、Q-Q Plot & 检查厚尾及偏度 \\
Temporal Dependence & ACF of Returns & 检查线性依赖 \\
Volatility Dependence & ACF of Absolute / Squared Returns & 检查波动率聚集 \\
Rolling Statistics & Rolling Mean / Variance & 调查潜在结构变化 \\
Information Availability & Timestamp Audit & 排除 Look-Ahead Bias \\
Sample Size Analysis & Effective Sample Size & 判断评估不确定性 \\
\bottomrule
\end{longtable}
\end{readertable}

一个最简单的 Python 实现可以写成：

\begin{verbatim}
import numpy as np
import pandas as pd
import matplotlib.pyplot as plt

from statsmodels.graphics.tsaplots import plot_acf
from scipy.stats import probplot

# Assume df contains chronologically ordered prices.
df = df.sort_index()

r = np.log(df["close"]).diff().dropna()

fig, axes = plt.subplots(2, 3, figsize=(14, 7))

# Distribution
axes[0, 0].hist(r, bins=60, density=True)
axes[0, 0].set_title("Return Distribution")

probplot(r, dist="norm", plot=axes[0, 1])
axes[0, 1].set_title("Normal Q-Q Plot")

# Rolling statistics
r.rolling(60).std().plot(ax=axes[0, 2])
axes[0, 2].set_title("Rolling Volatility")

# Temporal dependence
plot_acf(r, lags=30, ax=axes[1, 0])
plot_acf(r.abs(), lags=30, ax=axes[1, 1])
plot_acf(r**2, lags=30, ax=axes[1, 2])

axes[1, 0].set_title("ACF of Returns")
axes[1, 1].set_title("ACF of Absolute Returns")
axes[1, 2].set_title("ACF of Squared Returns")

plt.tight_layout()
plt.show()
\end{verbatim}

这只是探索性数据分析的参考代码，尚未在课题组数据上运行。实际使用时需要考虑复权价格、缺失值、交易日对齐及 ACF 置信区间的适用条件。

尤其需要注意，如果这些探索性分析用于设计 RSI 策略，那么应当在开发数据上进行，而不是提前检查封闭的最终测试集。

\hypertarget{part-4-52-0680-literature-index}{%
\subsection{11. Literature Index}\label{part-4-52-0680-literature-index}}

本阶段只推荐两项主要阅读材料，没有必要立刻进入完整的 Financial Econometrics 教材。

\textbf{P0 --- Empirical Properties of Asset Returns: Stylized Facts and Statistical Issues}

Rama Cont · 2001 · Quantitative Finance, 1(2), 223--236 · 同行评审

本阶段的核心参考文献。重点阅读收益率分布、非线性依赖、波动率聚集及统计估计困难。

\href{https://doi.org/10.1080/713665670}{DOI: 10.1080/713665670}

\textbf{P1 --- Revisiting Cont's Stylized Facts for Modern Stock Markets}

2025 · Quantitative Finance · 同行评审

重新检验经典 Stylized Facts 在现代市场中的表现，帮助我们避免将经验统计规律视为普遍成立的数学定律。

\href{https://doi.org/10.1080/14697688.2025.2530637}{DOI: 10.1080/14697688.2025.2530637}

第一篇的优先级明显更高。它能够为我们之后理解金融 Low-SNR Environment 提供统计学背景。

\hypertarget{part-4-52-0680-iv-a-conclusion-index}{%
\subsection{12. IV-A --- Conclusion Index}\label{part-4-52-0680-iv-a-conclusion-index}}

\hypertarget{part-4-52-0680-core-conclusions}{%
\subsubsection{12.1 Core Conclusions}\label{part-4-52-0680-core-conclusions}}

\begin{readertable}
\begin{longtable}[]{@{}
  >{\raggedright\arraybackslash}p{(\columnwidth - 4\tabcolsep) * \real{0.13}}
  >{\raggedright\arraybackslash}p{(\columnwidth - 4\tabcolsep) * \real{0.70}}
  >{\raggedright\arraybackslash}p{(\columnwidth - 4\tabcolsep) * \real{0.13}}@{}}
\toprule
ID & Key Conclusion & Index \\
\midrule
\endhead
A1 & Simple Return 通过乘法累积，Log Return 通过加法累积 & §1 \\
A2 & 一条历史序列只是随机过程的一条 Sample Path & §2.1 \\
A3 & Filtration 定义金融预测中的合法信息集合 & §2.2 \\
A4 & Conditional Mean 与 Conditional Variance 是不同预测目标 & §3 \\
A5 & Stationarity 不等于 Independence & §4 \\
A6 & Random Walk 的价格非平稳，但其增量可以平稳 & §4.3 \\
A7 & ACF 衡量时间序列的线性依赖 & §5 \\
A8 & 零自相关不意味着不存在非线性时间依赖 & §6.3 \\
A9 & Volatility Clustering 不必然意味着 Non-stationarity & §6.4 \\
A10 & Heavy Tails 可能限制传统统计推断方法的适用性 & §7 \\
A11 & Temporal Dependence 会改变有效样本信息量 & §8 \\
A12 & 金融数据特性可能降低 RSI Evaluator 的可靠性 & §9 \\
\bottomrule
\end{longtable}
\end{readertable}

\hypertarget{part-4-52-0680-core-formula-sheet}{%
\subsubsection{12.2 Core Formula Sheet}\label{part-4-52-0680-core-formula-sheet}}

\textbf{Result 1 --- Return Transformation}

\[r_t^{\log}=\log(1+r_t).\]

\textbf{Result 2 --- Financial DGP}

\[r_{t+1}
=
\mu_t+\sigma_tZ_{t+1}.\]

\textbf{Result 3 --- AR(1)}

\[X_t=\phi X_{t-1}+\varepsilon_t,
\qquad
\rho_h=\phi^{|h|}.\]

\textbf{Result 4 --- GARCH(1,1)}

\[\sigma_t^2
=
\omega+a\varepsilon_{t-1}^2+b\sigma_{t-1}^2.\]

在相应平稳条件下：
\[\mathbb E[\sigma_t^2]
=
\frac{\omega}{1-a-b}.\]

\textbf{Result 5 --- Effective Sample Size for AR(1)}

\[n_{\mathrm{eff}}
\approx
n\frac{1-\phi}{1+\phi}.\]

以上公式应结合对应模型的假设使用，而不是直接作为所有金融数据的固有性质。

\hypertarget{part-4-52-0680-research-takeaway}{%
\subsection{13. Research Takeaway}\label{part-4-52-0680-research-takeaway}}

如果将本阶段压缩为一个核心观点：

\textbf{金融 RSI 所面临的困难，不仅是需要从高噪声数据中找到规律，还需要在同一批有限、相关且可能非平稳的数据中，反复判断哪些研究修改是真正有效的。}

从这个角度来看，我们可以区分两个问题：
\[\underbrace{
\text{Can the agent find a better method?}
}_{\text{Optimization Capability}}\]

与：
\[\underbrace{
\text{Can the evaluator recognize a better method?}
}_{\text{Statistical Reliability}}\]

前者涉及 Part III，后者涉及 Part II 和 Part IV。

当两者的发展速度不匹配时，一个越来越强大的 Agent 可能反而越来越容易受到评估噪声的误导。这是我们目前提出的候选退化机制之一，尚需通过实际实验验证。

\hypertarget{part-4-52-0680-next-iv-b-low-signal-to-noise-ratio-predictability}{%
\subsection{Next --- IV-B: Low Signal-to-Noise Ratio \& Predictability}\label{part-4-52-0680-next-iv-b-low-signal-to-noise-ratio-predictability}}

IV-A --- Financial Time Series \& DGP

IV-B --- Low Signal-to-Noise Ratio \& Predictability

IV-C --- Financial Factors \& ML Models

IV-D --- Returns, Sharpe Ratio \& Evaluation

IV-E --- Backtest Overfitting, PBO \& DSR

IV-F --- Market Regime \& Robustness

下一阶段是 Part IV 中理论优先级最高的一节。

\begin{readertable}
\begin{longtable}[]{@{}
  >{\raggedright\arraybackslash}p{(\columnwidth - 2\tabcolsep) * \real{0.41}}
  >{\raggedright\arraybackslash}p{(\columnwidth - 2\tabcolsep) * \real{0.55}}@{}}
\toprule
Topic & Core Question \\
\midrule
\endhead
Signal-to-Noise Ratio & 如何定量描述金融预测的困难？ \\
Law of Total Variance & 如何分解可预测与不可预测成分？ \\
Irreducible Error & 模型性能的理论上限在哪里？ \\
Bayes Optimal Predictor & 最好的预测器在数学上是什么？ \\
Population \(R^2\) & 低信噪比如何限制预测能力？ \\
Efficient Market Hypothesis & 为什么金融规律可能难以持续利用？ \\
Spurious Correlation & 为什么无效金融特征也可能获得极高得分？ \\
Noise-driven Optimization & RSI 如何逐步将噪声误认为真实改进？ \\
\bottomrule
\end{longtable}
\end{readertable}

尤其是下一阶段，我们将利用你已经掌握的 \textbf{Conditional Expectation 与 Law of Total Variance}，推导一个相当漂亮的数学关系：
\[\boxed{
R_{\mathrm{optimal}}^2
=
\frac{\mathrm{SNR}}{1+\mathrm{SNR}}
}\]

在平方损失、有限二阶矩及相应 SNR 定义下，这个等式将金融数据的信噪比与理论最优预测能力直接联系起来。

它也将帮助我们更深入地理解，为什么 Low-SNR Environment 是 X Program \#161 中一个如此重要的限定条件。
